\documentclass[10pt,a4paper]{amsart}
\usepackage[margin=19mm]{geometry}
\usepackage{amsmath,amssymb,mathtools}
\usepackage[scr=rsfs,cal=euler]{mathalfa}
\usepackage{xfrac}
\usepackage[unicode,hidelinks,bookmarks=false]{hyperref}
\newcommand{\ie}{i.e.\ }
\newcommand{\cf}{cf.\ }
\newcommand{\resp}{respectively\ }
\newcommand{\Subsection}{\subsection}
\providecommand{\nobreakdash}{\nobreak\hbox{-}}
\setlength{\emergencystretch}{2em}
\multlinegap0pt
\usepackage{bm}
\usepackage{bbm}
\usepackage{dsfont}
\usepackage{xspace}
\usepackage{cancel}
\usepackage{mathtools}
\usepackage{amsthm}
\makeatletter
\@ifundefined{thm}{\newtheorem{thm}{Theorem}}{}
\makeatother
\title[Signed Mahonian Polynomials on Colored Derangements]{Signed Mahonian Polynomials on Colored Derangements}
\author{Hasan Arslan, Moussa Ahmia, Nazmiye Alemdar}
\date{Source: arXiv:2512.02404v2 (CC BY 4.0)}
\begin{document}
\maketitle

\noindent\emph{Source and licence.} Signed Mahonian Polynomials on Colored Derangements. Version arXiv:2512.02404v2 (CC BY 4.0), distributed under the Creative Commons Attribution 4.0 International licence, as recorded in the arXiv licence metadata for this exact version and shown on the abstract page https://arxiv.org/abs/2512.02404v2. Full rights notice: licenses/arxiv-2512.02404-CC-BY-4.0.txt.\par\smallskip

\noindent\textit{Editorial scope.} Counted unit: the Thm whose source label is \texttt{thm:signed-strong-derangements} in the article (source0.tex lines 834--844, source label \texttt{thm:signed-strong-derangements}), delivered together with the article's own complete original proof of it (source0.tex lines 846--968, one \texttt{proof} environment of 123 lines). No mathematics is written, repaired or re-derived: statement and proof are the article's own text line for line. Required context retained uncounted: none. Every definition, display and label of those lines is retained unchanged. Omitted from this package and not redistributed: 1--833, 845--845, 969--1738. Nothing inside a retained excerpt is omitted or altered: every retained body is delivered exactly as the article's lines are written, so no omission receipt of any kind accompanies this package. Citations: the retained excerpts carry no citation command at all, so no bibliography entry of the article is transcribed and no reference is redistributed. Reference adaptation: none was needed and none was made; every reference inside the delivered unit resolves inside it, so no unresolved reference or citation is produced. Printed slips kept literal: no printed word, symbol, hypothesis or number of the counted statement or of its proof was altered. Layout and class: the article's own class is \texttt{article} with packages including geometry, latexsym, amssymb, inputenc, enumerate, enumitem, color, graphicx, rotating, multirow, cite, amsthm, caption, tikz; the wrapper is the lane's pinned \texttt{amsart} editorial wrapper at 10pt on A4, which restates the theorem-like environments the retained text needs and adds the article's own macro definitions that the retained text needs (none), with extra wrapper packages (bm, bbm, dsfont, xspace, cancel, mathtools). The delivered numbering is the unit's own. Assets: no retained span contains a figure environment, a graphics-inclusion command, a PDF-page inclusion, a table or a TikZ picture; no image file is copied into this package, the package declares no asset dependency and no image file is redistributed with it. Licence: version arXiv:2512.02404v2 of this article is distributed under the Creative Commons Attribution 4.0 International licence, as recorded in the arXiv licence metadata for this exact version and shown on the abstract page https://arxiv.org/abs/2512.02404v2. A full rights notice is delivered at licenses/arxiv-2512.02404-CC-BY-4.0.txt. Characters: every retained excerpt is pure ASCII, so the delivered engine, XeLaTeX with its default Latin Modern text font, reports no missing character.
\begin{thm}\label{thm:signed-strong-derangements}
The difference between the number of even and odd colored derangements (according to the length function $L$) in $G_{c,n}$ is:
\[
d_{n}^{+(c)}-d_{n}^{-(c)}=\sum_{\pi\in D_n^{(c)}} (-1)^{L(\pi)}
=
\begin{cases}
(-1)^{n}, & \text{if $c$ is even},\\[4pt]
(-1)^{n-1}(n-1), & \text{if $c$ is odd}.
\end{cases}
\]
\end{thm}

\begin{proof}
Let $\pi = (\sigma;c_1,\dots,c_n)$ be the colored permutation
\[\pi=\sigma_1^{[c_1]}\cdots \sigma_n^{[c_n]},\]
where $\sigma\in \mathfrak{S}_n$ and $c_j\in\{0,1,\dots,c-1\}$.  

\medskip
\begin{enumerate}
\item \textbf{The case $c$ even.}
Denote
\[
\operatorname{Fix}(\sigma)=\{i\in[n]\mid \sigma(i)=i\},
\qquad
k=|\operatorname{Fix}(\sigma)|.
\]

If $i\in\operatorname{Fix}(\sigma)$, the condition $\pi(i)\neq i^{[0]}$ forces $c_i\in\{1,\dots,c-1\}$;  
if $i\notin\operatorname{Fix}(\sigma)$, then $c_i\in\{0,\dots,c-1\}$.  
Consequently,
\begin{equation}\label{eq:color-sum}
\sum_{\substack{(c_1,\dots,c_n)\\ (\sigma;\mathbf{c})\in D_n^{(c)}}}
(-1)^{\sum_{i=1}^n c_i}
=
\Bigl(\sum_{a=0}^{c-1}(-1)^a\Bigr)^{\,n-k}
\Bigl(\sum_{a=1}^{c-1}(-1)^a\Bigr)^{\,k}. 
\end{equation}

From the definition of $L$, we have
\[
(-1)^{L(\pi)}=(-1)^{\operatorname{inv}(\sigma)}(-1)^{\sum_{i=1}^n c_i}.
\]

Hence
\[
\sum_{\pi\in D_n^{(c)}}(-1)^{L(\pi)}
=
\sum_{\sigma\in\mathfrak{S}_n}(-1)^{\operatorname{inv}(\sigma)}
\sum_{\substack{(c_1,\dots,c_n)\\ (\sigma;\mathbf{c})\in D_n^{(c)}}}
(-1)^{\sum_{i=1}^n c_i}.
\]

For even $c$,
\[
\sum_{a=0}^{c-1}(-1)^a=0,\qquad
\sum_{a=1}^{c-1}(-1)^a=-1.
\]

Thus the right-hand side of \eqref{eq:color-sum} vanishes unless $k=n$; i.e. $\sigma=\operatorname{id}$.  
For the identity permutation the inner sum equals $(-1)^n$, whence
\[
\sum_{\pi\in D_n^{(c)}}(-1)^{L(\pi)}
=(-1)^{\operatorname{inv}(\operatorname{id})}(-1)^n=(-1)^n.
\]

\item \textbf{The case $c$ odd.}

For $j\in[n]$ define
\[
N(\sigma,j):=\bigl|\{i<j\mid \sigma(i)<\sigma(j)\}\bigr|,
\]
the number of non-inversions of $\sigma$ to the left of position $j$.  
Then, from the definition of $L$, we obtain
\[
(-1)^{L(\pi)}=(-1)^{\operatorname{inv}(\sigma)
+\sum_{j:\,c_j\neq 0}N(\sigma,j)
+\sum_{j=1}^{n}c_j},
\]
and consequently
\begin{equation}\label{eq:L-factor}
(-1)^{L(\pi)}
=
(-1)^{\operatorname{inv}(\sigma)}
\prod_{j=1}^{n}(-1)^{c_j}
\prod_{j:\,c_j\neq 0}(-1)^{N(\sigma,j)}.
\end{equation}

Fix $\sigma\in \mathfrak{S}_n$ and sum \eqref{eq:L-factor} over all colorings
$(c_1,\dots,c_n)\in\{0,1,\dots,c-1\}^n$ satisfying the derangement condition
$c_j=0\Rightarrow\sigma(j)\neq j$.  
Define
\[
S(\sigma):=\sum_{\substack{c_1,\dots,c_n\\ c_j=0\Rightarrow\sigma(j)\neq j}}
(-1)^{L(\pi)}.
\]

Factor out $(-1)^{\operatorname{inv}(\sigma)}$ and, for each $j$, set $p_j=(-1)^{N(\sigma,j)}$.

\begin{itemize}
\item \textit{If $\sigma(j)=j$ (a fixed point)}, then $c_j\neq0$ and $c_j\in\{1,\dots,c-1\}$.  
The contribution from such a position is
\[
\sum_{c_j=1}^{c-1}(-1)^{c_j} p_j
=p_j\sum_{k=1}^{c-1}(-1)^{k}.
\]
Because $c$ is odd, $c-1$ is even, and $\sum_{k=1}^{c-1}(-1)^k=0$.  
Hence if $\sigma$ has \textbf{any} fixed point, $S(\sigma)=0$.

\item \textit{If $\sigma(j)\neq j$ (non-fixed point)}, then $c_j\in\{0,\dots,c-1\}$.  
The contribution is
\[
\sum_{c_j=0}^{c-1}
\Bigl(\mathbf{1}_{c_j=0}
+\mathbf{1}_{c_j\neq0}\,(-1)^{c_j}p_j\Bigr)
=1+p_j\sum_{k=1}^{c-1}(-1)^k=1.
\]
\end{itemize}

Thus only classical \textbf{derangements} $\sigma$ contribute, and for them
\[
S(\sigma)=(-1)^{\operatorname{inv}(\sigma)}.
\]

\medskip

Summing over all permutations yields
\[
\sum_{\pi\in D_n^{(c)}}(-1)^{L(\pi)}
=\sum_{\sigma\in\mathfrak{D}_n}(-1)^{\operatorname{inv}(\sigma)}=(-1)^{\,n-1}(n-1),
\]
where $\mathfrak{D}_n$ denotes the set of ordinary derangements in $\mathfrak{S}_n$.
\end{enumerate}

Combining the two cases completes the proof.
\end{proof}

\end{document}
